\begin{table}[bt]
\small
%\color{blue}
\papertext{\vspace{-1mm}}
\centering
\scalebox{0.7}{
\begin{tabular}{|c|c|c|c|c|c|c|c|}
\hline
& \multicolumn{4}{c|}{\textbf{Question Types}} & \multicolumn{3}{c|}{\textbf{\# of Hops}}  \\ \hline
& \textbf{Look-up} & \textbf{Aggregation} & \textbf{Reasoning} & \textbf{Judge} & \textbf{1-Hop} & \textbf{2-Hops} & $\geq$\textbf{3-Hops} \\ \hline
\textbf{Qasper} & 54.1\% & 4\% & 24.4\% & 17.6\% & 72.3\% & 22.2\% & 5.6\% \\ \hline 
\textbf{NL\_DEV} & 75.2\% & 4\% & 18\% & 2.8\% & 42.3\% & 52.1\% & 5.6\% \\ \hline 
\textbf{HotpotQA} & 68.1\% & 1.8\% & 22.8\% & 7.4\% & 7.6\% & 57.3\% & 35.1\% \\ \hline
\textbf{\rthree{CUAD}} & \rthree{17.1\%} & \rthree{6.3\%} & \rthree{5.2\%} & \rthree{71.4\%} & \rthree{7.3\%} & \rthree{58.5\%} & \rthree{34.1\%} \\ \hline 
\textbf{\rthree{PubMedQA}} & \rthree{6\%} & \rthree{2\%} & \rthree{51\%} & \rthree{41\%} & \rthree{9.6\%} & \rthree{56.2\%} & \rthree{34.2\%} \\ \hline 
\end{tabular}}
\caption{\small Characteristics of Questions across Workloads.}
\papertext{\vspace{-4mm}}
\label{tab:query-workloads}
\end{table}


\begin{table*}[bt]
\small
\centering
\papertext{\vspace{-13pt}}
\scalebox{0.65}{
\begin{tabular}{|c|c|c|c|c|c|c|c|c|c|c|}
\hline
& & \multicolumn{3}{c|}{\textbf{Qasper} (8972 Tokens)} & \multicolumn{3}{c|}{\textbf{NL\_DEV} (7957 Tokens)} & \multicolumn{3}{c|}{\textbf{HotpotQA} (1233 Tokens)} \\ \hline 
 \textbf{Type} & \textbf{Strategy} & \textbf{Size Ratio (\%)} & \textbf{Cost Ratio} & \textbf{Latency} & \textbf{Size Ratio (\%)} & \textbf{Cost Ratio} &\textbf{Latency} & \textbf{Size Ratio (\%)} & \textbf{Cost Ratio} &\textbf{Latency} \\ 
\hline
\multirow{6}{*}{\textbf{Pruning Strategies}} & LLM\_bottom\_up & 9.7 & \textbf{0.2}  & \textbf{2} & 13.6 & 0.3 & 2 & 19.8 & \textbf{0.3} & 1.2 \\
&LLM\_top\_down & 12.6 & 0.9  & 3 & 15.3 & 1.2 & 2.9 & 22.3 & 1.3 &1.9\\
&LLM\_adaptive & 9.8 & \textbf{0.2}  & 2.1 & 13.5 & 0.3 & 2 & 20 & \textbf{0.3} & \textbf{1.1} \\ \cline{2-11}
&Embedding\_bottom\_up & 10.5 & 0.36 & 2.1 & 12.2 & \textbf{0.2}& \textbf{1.9} &  20.2 & \textbf{0.3} & 1.3\\
&Embedding\_top\_down & 12.4 & 1.1 & 3.1 & 14.8 & 1.1& 2.6 & 23.2 & 1.2 & 2\\
&Embedding\_adaptive & 10.6 & 0.35 & 2.1 & 12.1 & \textbf{0.2}& \textbf{1.9} &  20.3 & \textbf{0.3} & 1.2\\
\hline  \hline
\multirow{12}{*}{\textbf{\shortstack{Two-Phase Strategy: \\Pruning Strategies\\ + Refinement Strategies}}} 
&LLM\_bottom\_up\_EXP & 3.8 & \textbf{1.2} & 8.4  & 4.8 & \textbf{1.1}& \textbf{8.5} & 7.7 & \textbf{1.2} & \textbf{3.5} \\
&LLM\_bottom\_up\_SEQ & 3.6 & \textbf{1.3} & 9.7 & 4.9 & 1.3& 12.2 & 7.8 & 1.5 & 3.8\\
&LLM\_top\_down\_EXP& 3.7 & 1.7 & 11.2 &  5.1 & 1.9& 10.5 & 8.1 & 2.2 & 4.1\\
&LLM\_top\_down\_SEQ& 4.1 & 2.1 & 13.1 & 4.9 & 2.4& 14.2 & 8 & 2.9 & 4.7\\
&LLM\_adaptive\_EXP & 3.7 & \textbf{1.2} & \textbf{8.1}  & 4.8 & \textbf{1.1}& \textbf{8.5} & 7.7 & \textbf{1.1} & \textbf{3.3}  \\ 
&LLM\_adaptive\_SEQ & 3.6 & \textbf{1.3} & 9.3 & 4.8 & \textbf{1.2}& 11.4 & 7.8 & 1.4 & \textbf{3.7}\\ \cline{2-11}
&Embedding\_bottom\_up\_EXP& 3.4 & \textbf{1.3} & \textbf{8.2}    & 4.6 & \textbf{1.1}& \textbf{8.1} & 7.9 & \textbf{1.2} & 4.1 \\
&Embedding\_bottom\_up\_SEQ& 3.4 & 1.4 & 11.3 & 4.7 & \textbf{1.2}& 9.9 & 7.8 & 1.4 & 4.5\\
&Embedding\_top\_down\_EXP& 4.3 & 1.9 & 12.9 & 5.2 & 2.2& 10.8&8.2 & 2.4 & 5.8\\
&Embedding\_top\_down\_SEQ& 4.1 & 2.4 & 15.2 &  4.8 & 2.7& 13.5& 7.8 & 2.9 & 6.4\\
&Embedding\_adaptive\_EXP& 3.3 & \textbf{1.3} & \textbf{8}    & 4.6 & \textbf{1.1}& \textbf{8.3} & 7.9 & \textbf{1.1} & 4 \\
&Embedding\_adaptive\_SEQ& 3.4 & 1.4 & 10.8 & 4.7 & \textbf{1.2}& 9.6 & 7.8 & \textbf{1.3} & 4.2\\
\hline
\end{tabular}}
\caption{\small Size Ratio, Cost Ratio, and Latency of Strategies in \code{Qasper}, \code{NL\_DEV}, and \code{HotpotQA} Workloads. The top-1 pruning strategies, and the top-3 two-phase strategies are highlighted in bold for both cost ratio and latency. \code{EXP} and \code{SEQ} stand for \code{Exponential-Greedy} and \code{Sequential-Greedy}. }
\papertext{\vspace{-2.5em}}
\techreport{\vspace{-2em}}
\label{tab:top-1}
\end{table*}




%LLM\_top\_down & 17 & 0.7  & 3 & 21.4 & 0.7& 3 & 29.3 & 0.6 &1.4\\
%LLM\_top\_down\_EXP& 3.7 & 2.8 & 12.9 &  5.1 & 3.8& 10.5 & 8.1 & 2.2 & 4.4\\
%LLM\_top\_down\_SEQ& 4.1 & 3.7 & 15.5 & 4.9 & 5.7& 15 & 8 & 2.9 & 4.9\\
% Embedding\_top\_down & 20.9 & 0.5 & 2.5 & 19.4 & 0.9& 2.6 & 28.0 & 0.8 & 2.1\\
% Embedding\_top\_down\_EXP& 4.3 & 3.2 & 12.9 & 5.2 & 3.6& 10.8&8.2 & 2.1 & 5.8\\
% Embedding\_top\_down\_SEQ& 4.1 & 4.5 & 19.2 &  4.8 & 4.9& 15.5& 7.8 & 2.9 & 6.4\\
% Divide\_and\_Conquer & 14 & 1.2 & 3.2 & 16.0  & 1.1& 3 & 22.4 & 1.2 & 1.6\\
% Divide\_and\_Conquer\_EXP& 3.7 & 3.1 & 12.3 & 4.7 & 2.4& 10.3 & 7.6 & 2.0 & 4.3\\
% Divide\_and\_Conquer\_SEQ& 3.6 & 4.1 & 15.9 & 4.7 & 3.1& 12.6 & 7.6 & 2.3 & 4.5\\

%\vspace{-1mm}
\section{Evaluation} 
\label{sec:exp}

\begin{table}[bt]
\small
%\color{blue}
\centering
\scalebox{0.75}{
\begin{tabular}{l|l|c|c|c|c|c}
\hline
\textbf{Method} & \textbf{Configuration} & \textbf{Qasper} & \textbf{NL\_DEV} & \textbf{HotpotQA} & \rthree{\textbf{CUAD}} & \rthree{\textbf{PubMedQA}} \\ \hline
\multirow{3}{*}{\textbf{LLM}} & \textbf{gpt-4o-mini} & 0.4 & 0.37 & 0.46  & \rthree{0.42} & \rthree{0.39}\\ 
&\textbf{gpt-4o} & 0.36 & 0.26 & 0.3 & \rthree{0.31} & \rthree{0.34} \\ 
&\textbf{gemini-2-flash} & 0.6 & 0.65 & 0.47 & \rthree{0.47} & \rthree{0.42}\\ \hline 
\multirow{3}{*}{\textbf{Retrieval}} & \textbf{RAG-1\%} & 0.82 & 0.84  & 0.89  & \rthree{0.78} & \rthree{0.92}\\ 
&\textbf{RAG-5\%} & 0.59 & 0.64 & 0.7  & \rthree{0.66} & \rthree{0.81}\\ 
&\textbf{RAG-10\%} & 0.32 & 0.41 & 0.54 & \rthree{0.48} & \rthree{0.61}\\ \hline 
\textbf{BLIP}& {\bf E\_adaptive\_EXP} & \textbf{0.033} & \textbf{0.046} & \textbf{0.079} & \rthree{\textbf{0.03}} & \rthree{\textbf{0.12}}\\ \hline 
\end{tabular}}
\caption{\small \revised{Human Review Effort. (E stands for Embedding in BLIP.)}} 
\label{tab:human_review_time} 
\papertext{\vspace{-9mm}}
\techreport{\vspace{-7mm}}
\end{table}



We evaluate \trs{the performance of} \sys 
through experiments conducted on \rthree{five real-world}, two synthetic workloads, \rmix{as well as a user study}. 
%\agp{FIX! No longer three real-world workloads}
%\agp{Also mention user study}


% \begin{table}[bt]
% \small
% \color{orange}
% \centering
% \scalebox{1}{
% \begin{tabular}{c|c|c|c}
% \hline
% \textbf{Method}  & \textbf{Size Ratio(\%)} & \textbf{Cost Ratio} & \textbf{Latency}  \\ \hline
% LLM\_adaptive       &                  11.8 &  0.27 &  2.4 \\ \hline 
% Embedding\_adaptive   &        11.3 &  0.29 &  2.6 \\ \hline 
% LLM\_adaptive\_EXP   &            2.9 &    1.4 &    9.3 \\ \hline
% LLM\_adaptive\_SEQ   &            3.1 &    1.3 &    10.7 \\ \hline
% Embedding\_adaptive\_EXP &  2.8 &    1.4 &   9.8 \\ \hline
% Embedding\_adaptive\_SEQ &    2.9 &    1.5 &   11.2 \\ \hline 
% \end{tabular}}
% \caption{\small \rthree{Evaluation Results on CUAD and PubMedQA Workloads.}} 
% \label{tab:cuad} 
% \vspace{-1em}
% \end{table}

\begin{table}[bt]
\scriptsize
%\papertext{\color{orange}}
\centering
\scalebox{1}{
\begin{tabular}{c|c|c|c}
\hline
\textbf{Dataset}  & \textbf{Size Ratio(\%)} & \textbf{Cost Ratio} & \textbf{Latency}  \\ \hline
CUAD &  2.8 &    1.4 &   9.8 \\ \hline
PubMedQA & 11.3 & 1.2 & 4.2 \\ \hline 
\end{tabular}}
\caption{\small \rthree{Evaluation Results of \sys (the overall recommended strategy, Embedding\_adaptive\_EXP) on CUAD and PubMedQA.}} 
\label{tab:cuad} 
\papertext{\vspace{-3mm}}
\end{table}

\begin{figure*}[tb]
\centering
\subfigure[Ranker = Embedding\label{fig:seq-exp-embedding}]{\includegraphics[width=0.4\textwidth]{figures/seq_exp_embedding.pdf}}
\subfigure[Ranker = LLM\label{fig:seq-exp-llm}]{\includegraphics[width=0.4\textwidth]{figures/seq_exp_LLM.pdf}}
\vspace{-1.7em}
\caption{\small Cost ratios of pruning strategies on three real-world workloads and their skewed versions. }%\yiming{can remove figures of original workloads later since they are duplciated in Table 2}}  
\label{fig:cost-prune}
\vspace{-1.5em}
\end{figure*}





%\papertext{\vspace{-5mm}}
\subsection{Evaluation Setup}
\label{subsec:eval-setup}

\topic{Datasets \& Workloads} \rthree{We conduct experiments on five  commonly used question-answering workloads: {\em Qasper~\cite{dasigi2021dataset}}, {\em NL\_DEV}~\cite{kwiatkowski2019natural}, {\em HotpotQA~\cite{yang2018hotpotqa}}, {\em CUAD}~\cite{hendrycks2021cuad}, and {\em PubMedQA}~\cite{jin2019pubmedqa}, covering \techreport{the domains of} scientific papers, Wikipedia, legal contracts, and medical records.} For each workload, we sampled 500 distinct question-document pairs, ensuring that each question and document in a pair is unique. \rthree{The average tokens per document is 8,972, 7,957, 1,233, 9,964, and 891 for Qasper, NL\_DEV, HotpotQA, CUAD, and PubMedQA, respectively.} \revised{Questions with diverse types and complexities are evaluated: about \rthree{72\%} require at least two reasoning hops, as shown in Table~\ref{tab:query-workloads}; \papertext{see~\cite{blip} for details. }}\techreport{Here, Look-up and Aggregation denote factual questions, Judge corresponds to boolean questions (i.e., yes/no), and Reasoning questions ask why or how, typically requiring longer explanatory answers. } 
We additionally include two synthetic workloads, {\em Movie} and {\em Restaurant}, from TQA-bench~\cite{wolff2025well}, which focus on question answering over tables (TableQA), each containing  64,000 tokens and 100 questions.   

True provenance is available for TableQA workloads but not for others. Since the latter five workloads are on question answering over plain text, true provenance is unknown, and even human annotators may disagree on the ground truth; there may also be multiple true provenances. 
For the TableQA workloads, Movie and Restaurant, true provenance (i.e., the set of tuples) is provided for question-answer pairs. We focus on multi-hop reasoning questions in TableQA workloads, which are common analytical queries over tables. Out of 100 questions per TableQA workload, half are single-hop reasoning questions, while the others involve 2 to 3 hops.  

%According to ~\cite{wolff2025well}, LLMs have less than 10\% accuracy on  questions involving aggregates or joins over large tables (e.g., 64k tokens), limiting the availability of sufficient questions of this form for which  LLMs return correct answers. Morever, as we saw in Section~\ref{sec:database-prov-relationship}, provenance for queries that are not in $\mathcal{R}\mathcal{A}^+$ may not be useful. 
%We therefore focus on look-up questions, which are common analytical queries over tables, where LLMs can produce  accurate answers (e.g., "Which war movies from 2001 have a rating higher than 5.7?"). Out of the 100 questions per dataset, half are single-hop look-up questions (i.e., equivalent to SQL queries with one predicate), while the others involve 2 to 3 hops. 

%\revised{Our workloads}

\topic{Compared Strategies} We first compare the five strategies in the pruning phase, along with ten combined two-phase strategies. %, where one is  comprising a strategy from $\mathcal{S}_{\text{prune}}$ and one from  $\mathcal{S}_{\text{refine}}$. 
\code{Sequential-Greedy} and \code{Exponential-Greedy} are abbreviated as \code{SEQ} and \code{EXP}, respectively.  So \code{LLM\_bottom\_up\_EXP} represents a strategy that uses \code{LLM\_bottom\_up} and \code{Exponential-Greedy} in the pruning and refinement phase, respectively.  \revised{We also compare \sys against a retrieval-based approach, RAG, using various context sizes (1\%, 5\%, and 10\% of the full text), and against LLM-generated provenance as mentioned in Section~\ref{sec:introduction}.}  


\topic{Implementation} When performing \code{prune} in Algorithm~\ref{alg:linear-search}, we divide the text into $m=20$ equal-sized blocks. The embedding- and LLM-based rankers used to sort the blocks are implemented with the open-source embedding model \code{all-mpnet-base-v2} and LLM \code{Mistral 7B}. Both  embeddings and relevance scores are computed offline\techreport{for each document-question pair}. For question answering and provenance inference, we use \code{gpt-4o-mini} and \code{gemini-2-flash}. Results reported here are based on \code{gpt-4o-mini}. \code{Gemini-2-flash} has similar findings, \papertext{see \vldb{\cite{blip}}.}\techreport{in Table~\ref{tab:top-1-gemini}.} 


\topic{Metric} \revised{We first report the {\em accuracy}, defined in Section~\ref{sec:introduction}, as the fraction of document–question pairs for which the inferred provenance can reproduce an answer equivalent to the one from the full text.}  We further report the {\em size ratio}, defined as the number of tokens in the generated provenance divided by that in the full text. We additionally report the {\em latency} and the {\em cost ratio} that is defined as the cost of finding the provenance divided by the cost of question answering. 
For strategies using an embedding-based ranker, we exclude the latency for building embeddings and report only the latency for provenance inference, as embeddings can be computed offline. 
\revised{Finally, observing that an incorrect provenance requires humans to read the entire document, we report {\em human review effort}, measured as the fraction of the document corresponding to the provenance (i.e., size ratio) if accurate (i.e., the provenance is verifiable), and as one plus the size ratio otherwise.} 

We also measure whether methods can recover the true provenance $P_{true}$ for  TableQA workloads  by introducing two metrics: recovery ($R$) and exact-recovery ($eR$). 
\techreport{Provenance is represented as a set of tuples in the tables.} $R$ is defined as the percentage of questions where a predicted provenance $P_{pre}$ contains $P_{true}$ \techreport{over all questions}, i.e., $R = \frac{|P_{true} \subseteq P_{pre}|}{|QA|}$, while $eR$ requires $P_{pre}$ to be an exact match with $P_{true}$, i.e., $eR = \frac{|P_{true} = P_{pre}|}{|QA|}$. 
\rtwo{We note that the {\em accuracy} defined measures how often a provenance inferred by \sys can reproduce the answer (i.e., is verifiable), while {\em recovery} measures whether the inferred provenance matches the true provenance.} %\yiming{check the terminology of behavioral provenance is consistent with the new section 2.5}


%We measure the accuracy of provenance for the TableQA workloads and introduce two variants: exact accuracy ($eAccu$) and relaxed accuracy ($rAccu$). Provenance is represented as a set of tuples in the tables, and accuracy is defined as the percentage of correctly predicted provenance over all questions. $eAccu$ counts a correct provenance prediction as an exact match, i.e., $eAccu = \frac{|P_{true} = P_{pre}|}{|QA|}$, while $rAccu$ only requires the predicted provenance to be a superset of the ground truth provenance, i.e., $rAccu = \frac{|P_{true} \subseteq P_{pre}|}{|QA|}$. 









\vspace{-1em}
\subsection{Experiments}  
\label{subsec:exp-result}

\topic{Experiment 1: Cost, Latency, and Provenance Size of \sys}  
We present the cost ratio, latency, and the size ratio of the provenance generated by \sys in Table~\ref{tab:top-1}. We make several observations. First, strategies in the pruning phase effectively eliminate text portions unrelated to answers, with adaptive strategies  reducing average provenance size ratios to 10.2\%, 12.8\%, and 20.2\% for Qasper, NL\_DEV, and HotpotQA,  respectively. These strategies incur only 0.28$\times$, 0.25$\times$, and 0.3$\times$ the cost of answering questions over the full text, with latencies averaging 2.1, 1.9, and 1.1 seconds. Among them, \code{LLM\_adaptive} {\bf \em reduces provenance size to around 9.8\% of the full text, achieving a 0.2$\times$ cost ratio and an average latency of 2.1 seconds} across the workloads. 

Strategies in $\mathcal{S}_{refine}$ further reduce the provenance size ratio to about {\bf \em 3.7\%, 4.7\%, and 7.7\% on Qasper, NL\_DEV, and HotpotQA}, respectively. Here, \code{Exponential\_Greedy} (i.e., strategies ending with \code{EXP}) have up to a {\bf \em  1.1$\times$ cost ratio and an average latency of 8.1 seconds} across workloads, demonstrating that \delete{the two-phase strategies are  practical solutions for producing small-size provenances that reproduce the answer, while incurring only a slightly higher cost than answering the question over the full text.  Latency can be further improved via parallelization, as part of our future work}\cradd{two-phase strategies produce minimal provenances at only slightly higher cost than full-text QA. Latency can be further reduced via parallelization}. \rthree{We have consistent observations on the CUAD and PubMedQA  workloads, as shown in Table~\ref{tab:cuad}, where \sys incurs a slightly higher cost (with a cost ratio of 1.3) than question answering over the entire input to infer a minimal provenance efficiently. }











% \begin{table}[ht]
% \centering
% \small
% \scalebox{0.8}{
% \begin{tabular}{|c|c|c|c|c|c||c|c|}
% \hline
% & & \multicolumn{4}{c|}{{\bf Bottom-up (BU) VS Top-down (TD)}} & \multicolumn{2}{c|}{{\bf EXP VS SEQ}} \\ \hline 
%  \textbf{Dataset} & \textbf{Ranker}& \textbf{$CP_1$ (Percentile)} & \textbf{Q1} & \textbf{Median} & \textbf{Q3} & $\mathcal{S}_{Prune}$ & $CP_2$ \\
% \hline
% \multirow{4}{*}{\textbf{Qasper}} &  \multirow{2}{*}{LLM} & \multirow{2}{*}{7.9 (98.1\%)}  & \multirow{2}{*}{1} &\multirow{2}{*}{1} &\multirow{2}{*}{2} & BU & 0 \\
% & & & & & & TD & 0 \\ \cline{2-8}
%                     & \multirow{2}{*}{Embedding} & \multirow{2}{*}{7.1 (97.9\%)} & \multirow{2}{*}{1}& \multirow{2}{*}{1} &\multirow{2}{*}{3} & BU & 0 \\
% & & & & & & TD & 0 \\ \cline{1-8}
% %\hline
% \multirow{4}{*}{\textbf{NL\_DEV}} & \multirow{2}{*}{LLM} & \multirow{2}{*}{6.9 (95.6\%)} &\multirow{2}{*}{1} &\multirow{2}{*}{1} & \multirow{2}{*}{3} & BU & 0 \\
% & & & & & & TD & 0 \\ \cline{2-8}  
%                     & \multirow{2}{*}{Embedding} & \multirow{2}{*}{6.7 (93.6\%)}  &\multirow{2}{*}{1} & \multirow{2}{*}{1}& \multirow{2}{*}{2} & BU & 0 \\
% & & & & & & TD & 0 \\ \cline{1-8}
% %\hline
% \multirow{4}{*}{\textbf{HotpotQA}} & \multirow{2}{*}{LLM} & \multirow{2}{*}{7.2 (91.6\%)} &\multirow{2}{*}{1} & \multirow{2}{*}{1}& \multirow{2}{*}{3} & BU & 0\\
% & & & & & & TD & 0 \\ \cline{2-8}
%                     & \multirow{2}{*}{Embedding} & \multirow{2}{*}{7.4 (86.8\%)} &\multirow{2}{*}{1} &\multirow{2}{*}{1} & \multirow{2}{*}{2} & BU & 0 \\
% & & & & & & TD & 0 \\ %\cline{2-8}
% \hline
% \end{tabular}}
% \caption{\small Crossover Point (CP) and Statistics (Q1, Median, Q3) of Number of Returned Data Blocks.  }
% \label{tab:distribution}
% \end{table}


\vspace{1pt}
\topic{\revised{Experiment 2: Accuracy and Human Review Effort}}   \rthree{We have already presented the accuracy for retrieval-based (RAG) and LLM-generated provenance in Table~\ref{tab:agreement}, where \sys strategies guarantee {\em \bf an accuracy of 1, over 30\% higher than the best baseline} with a comparable provenance size (i.e., LLM-generated provenance using \texttt{gpt-4o}).} 
\rthree{\delete{Note that non-monotonic tasks account for around 5\%  across datasets (based on Table~\ref{tab:tasks}), and accuracy being 1 demonstrates that \sys is guaranteed to infer verifiable provenance. This is consistent with our theoretical result in Theorem~\ref{theo:optimal}.}\cradd{Accuracy of 1 on non-monotonic tasks ($\sim$5\% of queries; Table~\ref{tab:tasks}) is consistent with Theorem~\ref{theo:optimal}.} }

We report the human review effort in Table~\ref{tab:human_review_time}. \rthree{Compared with RAG-based and LLM-generated provenance, \sys requires substantially less human review effort---only {\bf \em 3\% to 12\%} across five workloads. } 
%{\bf \em 6.6\% to 7.3\% on Qasper, 9\% to 10.8\% on NL\_DEV, and 11\% to 19.3\% on HotpotQA}, respectively. 
This is because \sys always returns a minimal verifiable provenance, both effective and efficient for human verification.  


\techreport{
\begin{table}[]
\centering
\small
\scalebox{0.8}{
\begin{tabular}{|c|c|c|c|c|c||c|}
\hline
& & & \multicolumn{3}{c|}{{\bf Bottom-up vs. Top-down}} & {\bf EXP vs. SEQ} \\ \hline 
 \textbf{Dataset} & \textbf{Ranker}& {\bf $\mathcal{S}_{Prune}$} & \textbf{$CP_1$} & \textbf{Q1} & \textbf{Q3} &  {\bf $CP_2$ (Pct)} \\
\hline
\multirow{4}{*}{\textbf{Qasper}} &  \multirow{2}{*}{LLM} & BU & 8.2   & 1  &2  & 8.4 (35.2\%) \\
& & TD & 8.2  & 1 & 2  & 9.3 (32.7\%) \\ \cline{2-7}
                    & \multirow{2}{*}{Embedding} & BU & 7.4  & 1 &3  & 8.9 (31.9\%) \\
& & TD & 7.4  & 1 & 3 & 9.5 (28.6\%) \\ \cline{1-7}
%\hline
\multirow{4}{*}{\textbf{NL\_DEV}} & \multirow{2}{*}{LLM} & BU &6.9  &1 & 3 & 8.5 (38.4\%) \\
& & TD & 6.9  & 1 & 3 & 9.1  (35.2\%)\\ \cline{2-7}  
                    & \multirow{2}{*}{Embedding} & BU & 6.7   &1 &  2  & 8.2 (40.1\%) \\
& & TD & 6.7 & 1 & 2 &  8.9 (38.5\%) \\ \cline{1-7}
%\hline
\multirow{4}{*}{\textbf{HotpotQA}} & \multirow{2}{*}{LLM}&BU & 7.2  &1 & 3 & 7.5 (40.2\%)\\
& & TD & 7.2 &1 &3 & 8.1 (38.4\%) \\ \cline{2-7}
                    & \multirow{2}{*}{Embedding} &BU &  7.4  &1 & 2  & 7.2 (38.1\%) \\
& &TD & 7.4  & 1 & 2 & 8.9 (37.5\%) \\ %\cline{2-8}
\hline
\end{tabular}}
\caption{\small Crossover Points ($CP_1$ and $CP_2$). BU and TD stand for bottom\_up and top\_down, respectively. 
$Q1$ and $Q3$ are the first and third quartiles of the returned data blocks when comparing bottom-up and top-down. $Pct$ stands for the percentile of cases where the number of input sentences is less than $CP_2$.  }
\papertext{\vspace{-1.5em}}
\techreport{\vspace{-1em}}
\label{tab:distribution}
\end{table}
}

%\yiming{if no space, move table 9 to TR as it similarly examined pruning strategies on skewed workloads as figure 8. }


\techreport{\begin{table*}[bt]
\small
\centering
\scalebox{0.8}{
\begin{tabular}{|c|c|c|c|c|c|c|c|c|c|c|}
\hline
& & \multicolumn{3}{c|}{\textbf{Qasper} (8972 Tokens)} & \multicolumn{3}{c|}{\textbf{NL\_DEV} (7957 Tokens)} & \multicolumn{3}{c|}{\textbf{HotpotQA} (1233 Tokens)} \\ \hline 
 \textbf{Type} & \textbf{Strategy} & \textbf{Size Ratio (\%)} & \textbf{Cost Ratio} & \textbf{Latency} & \textbf{Size Ratio (\%)} & \textbf{Cost Ratio} &\textbf{Latency} & \textbf{Size Ratio (\%)} & \textbf{Cost Ratio} &\textbf{Latency} \\ 
\hline
\multirow{6}{*}{\textbf{Pruning Strategies}} & LLM\_bottom\_up & 10.3 & \textbf{0.3}  & 2 & 14.1 & 0.4 & 2.2 & 17.3 & 0.3 & 1.7 \\
&LLM\_top\_down & 13.4 & 0.8  & 2.7 & 14.7 & 1.4 & 3.2 & 19.5 & 1.1 &2.2\\
&LLM\_adaptive & 10.2 & \textbf{0.3}  & \textbf{1.8} & 13.9 & 0.3 & \textbf{2.1} & 17.7 & \textbf{0.2} & \textbf{1.6} \\ \cline{2-11}
&Embedding\_bottom\_up & 9.4 & 0.5 & 1.9 & 13.1 & \textbf{0.3}& 2.4 &  18.7 & 0.4 & 1.8\\
&Embedding\_top\_down & 10.8 & 1.2 & 2.8 & 14.5 & 1.2& 3.1 & 20.5 & 1.2 & 2.5\\
&Embedding\_adaptive & 9.7 & 0.4 & 2 & 13.1 & \textbf{0.3}& \textbf{2.2} &  18.9 & 0.4 & 1.7\\
\hline  \hline
\multirow{12}{*}{\textbf{\shortstack{Two-Phase Strategy: \\Pruning Strategies\\ + Refinement Strategies}}} 
&LLM\_bottom\_up\_EXP & 4.1 & \textbf{0.9} & \textbf{7.5}  & 5.2 & \textbf{0.9}& \textbf{9.1} & 7.3 & \textbf{0.8} & \textbf{3.7} \\
&LLM\_bottom\_up\_SEQ & 3.9 & \textbf{1.1} & 9.4 & 5 & \textbf{1.1}& 11.9 & 7.2 & 1.1 & 3.9\\
&LLM\_top\_down\_EXP& 4 & 1.4 & 10.4 &  5.1 & 1.6& 10.9 & 7.6 & 1.9 & 4\\
&LLM\_top\_down\_SEQ& 4.3 & 1.8 & 14.2 & 5.4 & 2.1& 13.7 & 7.8 & 2.3 & 4.3\\
&LLM\_adaptive\_EXP & 4 & \textbf{0.9} & 8.6  & 5.2 & \textbf{1.1}& \textbf{8.9} & 7.4 & \textbf{0.8} & \textbf{3.5}  \\ 
&LLM\_adaptive\_SEQ & 3.9 & \textbf{0.9} & 8.8 & 4.9 & \textbf{1}& 12.9 & 7.4 & \textbf{0.9} & \textbf{3.7}\\ \cline{2-11}
&Embedding\_bottom\_up\_EXP& 3.6 & 1.3 & \textbf{7.6}    & 5 & \textbf{1}& \textbf{8.4} & 7.3 & \textbf{1} & 3.9 \\
&Embedding\_bottom\_up\_SEQ& 3.7 & \textbf{1.1} & 10.6 & 4.9 & \textbf{1.1}& 9.6 & 7.6 & 1.4 & 4.2\\
&Embedding\_top\_down\_EXP& 3.9 & 1.7 & 11.3 & 5.2 & 1.7& 12.1&7.9 & 1.8 & 4.6\\
&Embedding\_top\_down\_SEQ& 4 & 2 & 13.9 &  5.3 & 2.2& 14.7& 7.8 & 2.3 & 5.7\\
&Embedding\_adaptive\_EXP& 3.6 & \textbf{1.1} & \textbf{8.2}    & 4.9 & \textbf{0.9}& \textbf{8.2} & 7.5 & \textbf{0.9} & \textbf{3.6} \\
&Embedding\_adaptive\_SEQ& 3.7 & \textbf{1.2} & 10.3 & 4.8 & \textbf{1}& \textbf{9.1} & 7.6 & \textbf{1} & 3.8\\
\hline
\end{tabular}}
\caption{\small Gemini-2-flash: Size Ratio, Cost Ratio, and Latency of Strategies in \code{Qasper}, \code{NL\_DEV}, and \code{HotpotQA} Workloads. The top-1 pruning strategies, and the top-3 two-phase strategies are highlighted in bold for both cost ratio and latency. \code{EXP} and \code{SEQ} stand for \code{Exponential-Greedy} and \code{Sequential-Greedy}. }
\vspace{-2em}
\label{tab:top-1-gemini}
\end{table*}}

\vspace{1pt}
\topic{Experiment 3: Bottom\_up, Top\_down, and Adaptive  strategies} 
Although bottom\_up and adaptive strategies outperform their top\_down counterparts in Table~\ref{tab:top-1}, we provide a breakdown to understand why. Section~\ref{subsubsec:bottom-top} shows that the cost of these strategies depends on the number of returned data blocks, with a suggested crossover point of $CP \approx 8.4$ when the text is divided into 20 blocks. We create a skewed dataset by upsampling question-document pairs so that 80\% of the pairs have more returned data blocks than $CP$. Figure~\ref{fig:cost-prune} reports the cost ratios on both the original and skewed workloads. \delete{It shows that top-down strategies are preferred when the number of returned data blocks is large—being over 2$\times$ cheaper than their bottom-up counterparts on the skewed workloads, with adaptive strategies providing a middle ground throughout.}\cradd{Top-down is over 2$\times$ cheaper on the skewed workloads, while adaptive strategies provide a middle ground in both settings.}
\techreport{To understand the trend of cost relative to the number of returned data blocks, we further plot the cost ratio of  \code{LLM\_bottom\_up},  \code{LLM\_top\_down}, and \code{LLM\_adaptive} for various numbers of returned data blocks in Qasper in Figure~\ref{fig:adaptive}. 
The cost of the bottom\_up strategy grows much faster than that of top\_down—it is lower when fewer than 8.2 data blocks are returned but higher beyond that. This is because top\_down examines blocks exponentially, while bottom\_up proceeds linearly, resulting in a roughly quadratic cost ratio relative to the number of returned data blocks.}  
\techreport{We also report the distribution of returned data blocks in Table~\ref{tab:distribution}, where the actual crossover points, denoted by $CP_1$, are presented when comparing bottom\_up and top\_down. The results show that most question-document pairs have a small number of returned blocks, making bottom\_up more cost-effective on average, whereas top\_down is preferred in the skewed workloads. }

%\delete{one strategy is to estimate the number of returned data blocks by posing an LLM call on the first $CP$ blocks computed by Theorem~\ref{theo:top-down-bottom-up}. If it reproduces the answer, suggesting fewer than $CP$ blocks are needed, then bottom\_up is chosen; otherwise, top\_down is preferred.} 
\delete{Overall, for pruning, if an estimate of the number of returned data blocks is available from historical data, one can choose a strategy by comparing this estimate with $CP$. Otherwise, the adaptive strategy wins, as it combines the strengths of both, being efficient for small outputs while avoiding quadratic cost for large ones.}\cradd{Overall, we recommend the adaptive strategy when provenance size is unknown, as it combines both strategies' strengths while avoiding worst-case costs.}




% \begin{table}[ht]
% \centering
% \small
% \scalebox{0.84}{
% \begin{tabular}{|c|c|c|c|c|c||c|}
% \hline
% & & & \multicolumn{3}{c|}{{\bf Bottom-up VS Top-down}} & {\bf EXP VS SEQ} \\ \hline 
%  \textbf{Dataset} & \textbf{Ranker}& {\bf $\mathcal{S}_{Prune}$} & \textbf{$CP_1$ (Pct)} & \textbf{Q1} & \textbf{Q3} &  {\bf $CP_2$ (Pct)} \\
% \hline
% \multirow{4}{*}{\textbf{Qasper}} &  \multirow{2}{*}{LLM} & BU & 8.2 (98.1\%)  & 1  &2  & 8.4 (35.2\%) \\
% & & TD & 8.2 (96.2\%) & 1 & 2  & 9.3 (32.7\%) \\ \cline{2-7}
%                     & \multirow{2}{*}{Embedding} & BU & 7.4 (97.9\%) & 1 &3  & 8.9 (31.9\%) \\
% & & TD & 7.4 (96.1\%) & 1 & 3 & 9.5 (28.6\%) \\ \cline{1-7}
% %\hline
% \multirow{4}{*}{\textbf{NL\_DEV}} & \multirow{2}{*}{LLM} & BU &6.9 (95.6\%) &1 & 3 & 8.5 (38.4\%) \\
% & & TD & 6.9 (93.2\%) & 1 & 3 & 9.1  (35.2\%)\\ \cline{2-7}  
%                     & \multirow{2}{*}{Embedding} & BU & 6.7 (93.6\%)  &1 &  2  & 8.2 (40.1\%) \\
% & & TD & 6.7(92.8\%) & 1 & 2 &  8.9 (38.5\%) \\ \cline{1-7}
% %\hline
% \multirow{4}{*}{\textbf{HotpotQA}} & \multirow{2}{*}{LLM}&BU & 7.2 (91.6\%) &1 & 3 & 7.5 (40.2\%)\\
% & & TD & 7.2(90.4\%) &1 &3 & 8.1 (38.4\%) \\ \cline{2-7}
%                     & \multirow{2}{*}{Embedding} &BU &  7.4 (86.8\%) &1 & 2  & 7.2 (38.1\%) \\
% & &TD & 7.4 (84.2\%) & 1 & 2 & 8.9 (37.5\%) \\ %\cline{2-8}
% \hline
% \end{tabular}}
% \caption{\small Crossover Points ($CP_1$ and $CP_2$). BU and TD stand for bottom\_up and top\_down, respectively. $Pct$ stands for the percentile of cases where the number of returned data blocks is less than $CP_1$, or the number of input sentences is less than $CP_2$. 
% $Q1$ and $Q3$ are the first and third quartiles of the returned data blocks when comparing Bottom-up and Top-down. }
% \vspace{-2em}
% \label{tab:distribution}
% \end{table}

\techreport{
\begin{figure}[tb]
    \centering
    \includegraphics[width=0.9\linewidth]{figures/LLM_bottom_up_top_down_gpt4omini_paper.pdf}
    \vspace{-1em}
    \caption{\small Cost Ratio of top\_down and bottom\_up on Qasper.} 
    \papertext{\vspace{-0.7em}}
    \label{fig:adaptive}
\end{figure}
}


% \begin{figure}[tb]
%     \centering
%     \includegraphics[width=0.8\linewidth]{figures/seq_exp_compare.pdf}
%     \vspace{-1em}
%     \caption{\small Cost Ratio of \code{EXP} and \code{SEQ} on Qasper.} 
%     \vspace{-0.2em}
%     \label{fig:cost-exp-seq}
% \end{figure}

% \begin{figure}[tb]
%     \centering
%     \begin{minipage}[t]{0.475\linewidth}
%         \centering
%         \includegraphics[width=\linewidth]{figures/LLM_bottom_up_top_down_gpt4omini_paper.pdf}
%         \vspace{-2em}
%         \caption{\small Cost Ratio of \textit{top\_down} and \textit{bottom\_up} on Qasper.}
%         \vspace{-0.5em}
%         \label{fig:adaptive}
%     \end{minipage}
%     \hfill
%     \begin{minipage}[t]{0.485\linewidth}
%         \centering
%         \includegraphics[width=\linewidth]{figures/seq_exp_compare.pdf}
%         \vspace{-2em}
%         \caption{\small Cost Ratio of \code{EXP} and \code{SEQ} on Qasper.}
%         \vspace{-0.5em}
%         \label{fig:cost-exp-seq}
%     \end{minipage}
% \end{figure}


%Table~\ref{tab:distribution} presents the distribution of returned data block numbers, including the percentage of question-document pairs (out of 500) with returned data blocks no greater than the cross point (CP), as well as the values for Q1 (25\%), Median, and Q3 (75\%). For most question-document pairs, the number of returned data blocks in bottom\_up strategies is small (less than 3, as indicated by Q3). Consequently, the average cost is dominated by bottom\_up strategies. The adaptive strategy leverages this advantage while enhancing robustness by avoiding high costs in cases with a large number of returned data blocks.

% \begin{table}[bt]
% \small
% \centering
% \scalebox{1}{
% \begin{tabular}{|c|c|c|c|c|}
% \hline
% & \multicolumn{2}{c|}{\textbf{Movie} (64k Tokens)} & \multicolumn{2}{c|}{\textbf{Restaurant} (64k Tokens)}  \\ \hline 
% & $eAccu$ & $rAccu$ &  $eAccu$ & $rAccu$ \\ \hline
% $Q_{simple}$  & 1 & 1 & 0.98 & 1  \\ \hline
% $Q_{complex}$ & 0.98 & 1 & 0.96 & 1 \\ \hline 
% \end{tabular}}
% \caption{\small Accuracy in TableQA Workloads}
% \label{tab:accuracy}
% \vspace{-1em}
% \end{table}




%Regardless of whether relevance is estimated using embedding similarities or LLMs, bottom-up strategies (e.g., \code{LLM\_bottom\_up}) consistently outperform their corresponding top-down strategies (e.g., \code{LLM\_top\_down}). In particular, \code{LLM\_bottom\_up} and \code{Embedding\_bottom\_up} have 62\% and 60\%  smaller provenance sizes than top-down counterparts, \code{LLM\_top\_down} and \code{Embedding\_top\_down}. For two-phase strategies when using an LLM for relevance estimation,  \code{LLM\_bottom\_up\_EXP} and \code{LLM\_bottom\_up\_SEQ} are  2.1$\times$ and 2.5$\times$ cheaper compared to their top-down counterparts, \code{LLM\_top\_down\_EXP} and \code{LLM\_top\_down\_SEQ}. When using embeddings, a similar trend is observed, with bottom-up strategies achieving 2.2$\times$ and 2.5$\times$ lower cost  than their top-down equivalents. 

%Out of the 20 data blocks in the given text $T$, the average number of returned blocks for the bottom\_up strategy is 2.1, 2.6, and 3.7 for Qasper, NL\_DEV, and HotpotQA, respectively. In contrast, the top\_down strategy returns 3.6, 3.7, and 5.8 blocks on average for the same workloads. 


%This is consistent with the theoretical analysis in Section~\ref{subsubsec:bottom-top}, where the actual number of blocks is smaller than $\sqrt{20}$, bottom\_up strategies are expected to perform better than top-down strategies. The large number of blocks returned by top-down also demonstrates that LLM performs less accurately on long contexts than short one.  In top-down strategies, LLM verification starts from large sequence, so a single false negative can cause early termination, leading to an overly large provenance. In contrast, bottom-up strategies begin with shorter sequences, where LLMs perform more reliably, making it easier to identify accurate top-$k$ blocks.

% \begin{figure*}[tb]
% \centering
% \vspace{-0.5em}
% \captionsetup[subfigure]{skip=0pt} % 紧凑子图标题间距
% \subfigure[Cost Ratio\label{fig:topk-cost}]{\includegraphics[width=0.32\textwidth]{figures/cost.pdf}}
% \subfigure[Latency\label{fig:topk-latency}]{\includegraphics[width=0.32\textwidth]{figures/latency.pdf}}
% \subfigure[Provenance Size\label{fig:topk-size}]{\includegraphics[width=0.32\textwidth]{figures/size.pdf}}
% \vspace{-1.5em}
% \caption{\small Cost Ratio, Latency, and Provenance Size Ratio of Top-$k$ Provenance.}
% \label{fig:top-k-eval}
% \vspace{-1.2em}
% \end{figure*}

% \begin{figure*}[tb]
% \centering
% \vspace{-0.5em}

% % Disable subfigure labels (for subcaption package)
% \captionsetup[subfigure]{labelformat=empty, skip=0pt}

% % No subcaptions and no (a)/(b)/(c)
% \subfigure{\includegraphics[width=0.32\textwidth]{figures/cost.pdf}}
% \subfigure{\includegraphics[width=0.32\textwidth]{figures/latency.pdf}}
% \subfigure{\includegraphics[width=0.32\textwidth]{figures/size.pdf}}
% \vspace{-2.2em}
% \caption{\small Cost Ratio, Latency, and Provenance Size Ratio of Top-$k$ Provenance.}
% \label{fig:top-k-eval}
% \vspace{-1.2em}
% \end{figure*}

% Requires: \usepackage{caption}
\begin{figure*}[tb]
    \vspace{-2mm}
\centering

\begin{minipage}[b]{0.24\textwidth}
    \centering
    \vspace{-2mm}
    \includegraphics[width=\linewidth]{figures/seq_exp_compare.pdf}
    \vspace{-8.3mm}
    \captionof{figure}{\small Cost Ratio. (\code{EXP} vs. \code{SEQ})}
    \label{fig:cost-exp-seq}
\end{minipage}
\hfill
\begin{minipage}[b]{0.75\textwidth}
    \centering
    \vspace{-2mm}
    \includegraphics[width=0.32\linewidth]{figures/cost.pdf}\hfill
    \includegraphics[width=0.32\linewidth]{figures/latency.pdf}\hfill
    \includegraphics[width=0.32\linewidth]{figures/size.pdf}
    \vspace{-5mm}
    \captionof{figure}{\small Cost Ratio, Latency, and Provenance Size Ratio of Top-$k$ Provenance.}
    \label{fig:top-k-eval}
\end{minipage}
\papertext{\vspace{-5mm}}
\end{figure*}


\begin{table}[bt]
\small
\vspace{-10pt}
\centering
%\color{blue}
\scalebox{0.63}{
\begin{tabular}{|c|c|c|c|c|c|c|c|c|c|c|c|c|c|}
\hline
\textbf{Methods} & \textbf{Models} & \multicolumn{6}{c|}{\textbf{Movie (64k Tokens)}} & \multicolumn{6}{c|}{\textbf{Restaurant (64k Tokens)}}  \\ \hline & & \multicolumn{3}{c|}{\textbf{$Q_{simple}$}} & \multicolumn{3}{c|}{\textbf{$Q_{complex}$}} & \multicolumn{3}{c|}{\textbf{$Q_{simple}$}} & \multicolumn{3}{c|}{\textbf{$Q_{complex}$}} \\ \hline 
& & \textbf{$eR$} & \textbf{$R$} & \rtwo{\textbf{$A$}} &  \textbf{$eR$} & \textbf{$R$} & \rtwo{\textbf{$A$}} & \textbf{$eR$} & \textbf{$R$} & \rtwo{\textbf{$A$}} & \textbf{$eR$} & \textbf{$R$} & \rtwo{\textbf{$A$}} \\ \hline
\textbf{\sys} & \textbf{E\_adaptive}  & \textbf{1} & \textbf{1} & \rtwo{\textbf{1}} & \textbf{0.98} & \textbf{1} & \rtwo{\textbf{1}} & \textbf{0.98} & \textbf{1}  & \rtwo{\textbf{1}} & \textbf{0.96} & \textbf{1} & \rtwo{\textbf{1}} \\ \hline
\multirow{3}{*}{\textbf{LLM}} & \textbf{gpt-4o-mini} & 0.62 & 0.76 & \rtwo{0.74} & 0.58 & 0.75 & \rtwo{0.73} & 0.72 & 0.84 & \rtwo{0.84} & 0.73 & 0.81 & \rtwo{0.79} \\  
&  \textbf{gpt-4o} & 0.64 & 0.89 & \rtwo{0.86} & 0.62 & 0.83 & \rtwo{0.82} & 0.74 & 0.9 & \rtwo{0.88} & 0.79 & 0.86 & \rtwo{0.86} \\ 
&  \textbf{gemini-2-flash} & 0.83 & 0.91 & \rtwo{0.9} & 0.67 & 0.8 & \rtwo{0.76} & 0.7 & 0.86 & \rtwo{0.83} & 0.73 & 0.79 & \rtwo{0.78} \\ \hline 
\multirow{4}{*}{\textbf{Retrieval}} & \textbf{RAG-EXACT} & 0.44 & 0.44 & \rtwo{0.44} & 0.32 & 0.32 & \rtwo{0.32} & 0.68 & 0.68 & \rtwo{0.68}  & 0.62 & 0.62 & \rtwo{0.62} \\ %\hline 
& \textbf{RAG-1\%} & 0 & 0.48 & \rtwo{0.48} &  0 & 0.42& \rtwo{0.41} & 0 & 0.82 & \rtwo{0.81}& 0 & 0.74& \rtwo{0.74} \\
%\hline 
& \textbf{RAG-5\%} & 0 & 0.72& \rtwo{0.7} & 0 & 0.58& \rtwo{0.56} & 0 & 0.92& \rtwo{0.91} & 0 & 0.86& \rtwo{0.86} \\ %\hline
& \textbf{RAG-10\%} & 0 & 0.8& \rtwo{0.76} & 0 & 0.74& \rtwo{0.71} & 0 & 0.96& \rtwo{0.93} & 0 & 0.94& \rtwo{0.93} \\ \hline
\end{tabular}}
\caption{\small \revised{Recovery} and \rtwo{Accuracy} in TableQA Workloads. (E stands for Embedding.)}
\label{tab:accuracy}
\vspace{-1em}
\end{table}

% %%% Temporary table: average accuracy (A) per dataset, averaged over Q_simple and Q_complex.
% \begin{table}[bt]
% \small
% \vspace{-10pt}
% \centering
% \scalebox{0.8}{
% \begin{tabular}{c|c|c|c}
% \hline
% \textbf{Methods} & \textbf{Models} & \textbf{Movie (64k Tokens)} & \textbf{Restaurant (64k Tokens)} \\ \hline
% \textbf{\sys} & \textbf{E\_adaptive} & \textbf{1} & \textbf{1} \\ \hline
% \multirow{3}{*}{\textbf{LLM}} & \textbf{gpt-4o-mini} & 0.74 & 0.82 \\
% & \textbf{gpt-4o} & 0.84 & 0.87 \\
% & \textbf{gemini-2-flash} & 0.83 & 0.81 \\ \hline
% \multirow{4}{*}{\textbf{Retrieval}} & \textbf{RAG-EXACT} & 0.38 & 0.65 \\
% & \textbf{RAG-1\%} & 0.45 & 0.78 \\
% & \textbf{RAG-5\%} & 0.63 & 0.89 \\
% & \textbf{RAG-10\%} & 0.74 & 0.93 \\ \hline
% \end{tabular}}
% \caption{\small Average \rtwo{Accuracy} ($A$) in TableQA Workloads, averaged over $Q_{simple}$ and $Q_{complex}$. (E stands for Embedding.)}
% \label{tab:accuracy-avg}
% \vspace{-1em}
% \end{table}


%1.48$\times$, 1.35$\times$, and 1.25$\times$ of cost,
\vspace{1pt}
\topic{Experiment 4: Exponential-Greedy (EXP) versus Sequential-Greedy (SEQ) Strategies} Consider \code{SEQ} and \code{EXP} \trs{in the refine phase} in Table~\ref{tab:top-1}.  \code{EXP} consistently outperforms \code{SEQ}, 11.6\%, 12\%, and 18.6\% cheaper, 
 and 1.3$\times$, 1.34$\times$, and 1.1$\times$ faster for Qasper, NL\_DEV, and HotpotQA, respectively. When the size of the input provenance is larger, \code{EXP} shows greater improvement over \code{SEQ}, up to 24\% cheaper and 1.49$\times$ faster. \delete{To provide a breakdown of their performance, we plot their cost ratios conditioned on the number of sentences in the provenance returned from the pruning phase in Figure~\ref{fig:cost-exp-seq} on Qasper. Both strategies perform similarly on small inputs, with \code{SEQ} slightly preferred when the number of sentences is less than 10\trs{ in Figure~\ref{fig:cost-exp-seq}}, while \code{EXP} scales better to larger inputs.}\cradd{Figure~\ref{fig:cost-exp-seq} plots their cost ratios vs.\ input size on Qasper: \code{SEQ} is preferred below 10 sentences, while \code{EXP} scales better beyond that.}   \techreport{Similar observations hold for other strategies and workloads in Table~\ref{tab:distribution}, where $CP_2$ denotes the crossover point based on the number of input sentences. Specifically, when the input contains fewer than $CP_2$ sentences, \code{SEQ} has a lower cost for more than half of the question-document pairs. An average percentile (Pct for $CP_2$) below 40\%, implying that only less than 40\% of inputs favor \code{SEQ}, explains why \code{EXP} has an overall lower cost. } 
Overall, when the number of sentences is small (10 is an empirical crossover point), \code{SEQ} is recommended for its simplicity and lower cost. Otherwise, \code{EXP} is clearly preferred. 


%This observation is consistent with the theoretical analysis presented in Section~\ref{subsubsec:seq-exp}. 
%Based on the above observations, we recommend {\bf \em combining \code{adaptive} with \code{Exponential-Greedy} to infer a minimal provenance}. For ranker models, while embeddings and LLMs do not show a clear advantage in all cases, embeddings are preferable when pre-built embeddings are available. Otherwise, LLMs can be explored as an alternative. 









\vspace{1pt}
\topic{Experiment 5: \revised{Recovery} of True Provenance in TableQA Workloads} 
To measure how well methods can recover the true provenance, we report both recovery ($R$) and exact recovery ($eR$) \rtwo{along with accuracy ($A$) for \sys} against four variants of RAG- and LLM-generated provenances in Table~\ref{tab:accuracy}.  \sys uses an adaptive strategy using embedding model as the ranker, where \code{SEQ} or \code{EXP} is chosen based on the provenance size (fewer than 10 sentences or not) returned in the previous phase. %\papertext{All RAGs use the same embedding model as \sys.} 
\techreport{All RAGs rank tuples based on embedding similarity to the question–answer pair, using the same embedding model as \sys.}  The top-$k$ tuples are used in RAG-EXACT, where $k$ matches the size of the true provenance. In the RAG-$p$\% strategy, the top-$p$\% of tuples are returned.
%as the predicted provenance. 

\sys can always recover the superset of true provenance as indicated by $R=1$, while it recovers the exact true tuples for over 98\% and 96\% of questions in Movie and Restaurant workloads, respectively. %demonstrating that \sys always returns the true provenance, and only returns a superset of the provenance in 2\%---4\% of the time. 
In contrast, RAG-EXACT achieves significantly lower $eR$, recovering the  true provenance in only 32\% to 68\% of cases across the two workloads.  Increasing the size of the returned provenance to 10\% (RAG-10\%) improves recovery but introduces many false positives (over 50 times larger than the ground-truth provenance), making it less interpretable. \techreport{For example, in Restaurant, RAG-10\% returns 106 tuples, while the true provenance contains only 1 tuple, undermining its usefulness.} 
LLM-based baselines recover the exact ground-truth provenance in 50--70\% of cases across different LLMs and workloads, and fail to capture the true in over 20\% of complex workloads.  \rtwo{Finally, the provenance inferred by \sys can always reproduce the answer, as indicated by $A=1$, while other approaches have no guarantee of producing verifiable provenance. }

%Considering $eAccu$, \sys recovers the exact ground truth  tuples for over 98\% and 96\% of questions in Movie and Restaurant workloads, respectively. Moreover, all predicted provenances contain the true provenance, as indicated by $rAccu = 1$, demonstrating that \sys accurately infers the correct provenance. 

\vspace{1pt}
\topic{Experiment 6: Top-k Provenance}  
%We begin by examining the number of returned distinct minimal provenances by running  \\ \code{top-k-provenance} with $k$ set to $+\infty$. Figure~\ref{fig:distribution} presents the first quartile (Q1, 25\% percentile), median, and third quartile (Q3, 75\% percentile) of the number of provenances identified across the Qasper, HotpotQA, and NL\_DEV workloads. As shown, Q1 and the median are 1 for all workloads, while Q3 ranges from 2 to 3. This suggests that for most document-question pairs across these workloads, the number of distinct minimal provenances is small, typically fewer than 3. 
We examine the performance of \code{top-k-provenance} with $k$ ranging from 1 to 5 in Figure~\ref{fig:top-k-eval}. 
 The algorithm stops until {\em up to} $k$ minimal provenances are found. %Figure~\ref{fig:top-k-eval} presents the cost ratio, latency, and provenance size ratio. 
As $k$ increases, the cost ratio (and similarly, latency) shows a modest increase, rising from approximately 1.3 to 1.8 across all workloads. This is because during the tree-based search in \code{top-k-provenance}, intermediate results computed for earlier provenances can often be reused when identifying subsequent ones. Consequently, later provenances frequently begin from smaller initial text sequences, leading to lower computation cost and latency. \delete{For example, identifying the provenance for node 9 in Figure~\ref{fig:top-k} starts from the text corresponding to node 9 rather than the full text, making the computation faster when returning the second provenance.}\cradd{For example, the second provenance starts from node 9's subtree in Figure~\ref{fig:top-k} rather than the full text, reducing cost.}  
\techreport{Morever, the average size of the returned minimal provenances remains stable, which is consistent with our findings in Table~\ref{tab:top-1}. }

% \subsection{User Studies}
% \label{subsec:user-study} 



\begin{table}[bt]
\papertext{\vspace{-10pt}}
\scriptsize
%\papertext{\color{myForestGreen}}
\centering
\scalebox{0.9}{
\begin{tabular}{|c|c|c|c|c|c|c|c|c|}
\hline
& \multicolumn{2}{c|}{\textbf{police\_reports}} & \multicolumn{2}{c|}{\textbf{legal}}  & \multicolumn{2}{c|}{\textbf{paper}} & \multicolumn{2}{c|}{\textbf{finance}} \\ \hline
& $H_q$ & $H_e$ & $H_q$ & $H_e$ & $H_q$ & $H_e$ & $H_q$ & $H_e$ \\ \hline 
\textbf{\sys} & \textbf{4.1}  & \textbf{75\%} &\textbf{5} & \textbf{100\%} &\textbf{4.9} & \textbf{88\%} &\textbf{5} & \textbf{63\%} \\ \hline 
\textbf{RAG} & 3.1 & 38\% &3.1 & 25\% &3.3 & 20\% &2.8 & 0\% \\ \hline 
\end{tabular}}
\caption{\rtwo{\small User Study. ($H_q$ ranges from 1 to 5 and measures the usefulness of the provenance, while $H_e$ is the percentage of cases in which a provenance contains the right amount of information.)}}
\papertext{\vspace{-3mm}}
\label{tab:user-study}
\end{table}

\vspace{-3mm}
\subsection{User Study}
\label{subsec:user-study}

\rtwo{Our experimental evaluation thus far
has focused on the accuracy and recoverability of provenance,
as opposed to its usefulness.
We conducted a user study to assess 
whether humans
find the provenance returned by \sys to be useful 
 when validating answers from LLMs.
We selected four documents from 
varied domains: %one each from 
{\em police\_reports}, 
{\em legal\_docs},  {\em scientific papers}, and {\em finance\_docs}, drawn from our legal partners~\cite{clean} and open-source benchmarks~\cite{dasigi2021dataset,finance}. 
We selected three questions per document, 
testing three scenarios: 
1) provenance generated by \sys; 
2) RAG-generated provenance, where up to 5\% (or 500 tokens, whichever is smaller) for ease of human interpretation; and 3) no provenance. 
Our goal was to compare the effectiveness of
\sys's provenance with that returned by RAG.
The last scenario was targeted at 
understanding human review effort to
determine correctness  
when no provenance is provided. 
We had eight participants in our user study,
who considered each scenario for each document, amounting
to 12 scenarios in total. 
\techreport{We randomized the order of settings across participants.}
\papertext{The details of other setups  
can be found in~\cite{blip}.}} 

\noindent \rtwo{{\bf Metrics.} 
Participants rated the {\em quality}
of a returned provenance from 1 (low)--5 (high),
assessing its usefulness in supporting
the answer returned by the LLM;
we let $H_q$ be the average quality across
participants for each dataset.
We also asked whether the provenance was both necessary and sufficient (i.e., containing no extra information and missing no critical information); let $H_e$ denote the percentage of such cases, averaged across participants per dataset. }

\rtwo{Table~\ref{tab:user-study} 
reports the average $H_q$ and $H_e$ over eight participants.
Overall, participants found the provenance 
inferred by \sys to be useful for judging answer correctness, 
as indicated by the overall $H_q$ of 4.75 (out of 5)\papertext{.}  
\techreport{across all documents.} 
One participant (P3) commented that ``{\em this seems to contain 
the right amount of information, 
including proper context to figure out the answer.}'' 
\delete{Cases where participants were not satisfied were
primarily due to semantic ambiguity\papertext{,}
\techreport{in the concepts contained in the provenance, }
where the LLM may carry different interpretations from humans.
For example, to identify officers using force with an answer ``John Smith'', when interpreting the provenance ``{\em Use of force details. Involved officer: reporting officer -- John Smith}\footnote{Name anonymized for privacy.}'',
the term ``reporting officer'' causes
ambiguity as to whether John
actually used force, indicating more context
would have been helpful for assessing validity.}\cradd{Dissatisfaction arose mainly from semantic ambiguity between LLM and human interpretations. For example, the provenance ``{\em Use of force details. Involved officer: reporting officer -- John Smith}\footnote{Name anonymized for privacy.}'' left it unclear whether John actually used force, suggesting more context could help.}}

\rtwo{RAG-based provenance was deemed to be 
less effective, with an overall $H_q$ of 3.1. One participant (P2) complained: “{\em This looks more like a search result than provenance, and provenance should explain why this answer. }” 
Participants also found \sys's 
provenance to contain the right amount 
of information 82\% of the time, 61\% higher 
than that for RAG-based provenance. 
\delete{When participants felt that \sys
returned extra information, this was typically because
some of this information was not necessary given
the participants' domain knowledge. For example, to identify the number of employees at Amazon from its 10-K financial report with the answer ``29,239'', users can interpret the provenance ``{\em we employed 29,239 people...}'' using their knowledge that ``we'' refers to Amazon. In contrast, the additional provenance returned by \sys that includes the earlier mention of ``Amazon'' (prior to ``we'') is often deemed unnecessary for users, though it is required for LLMs to interpret the answer.}\cradd{When \sys returned extra information, it was typically unnecessary given users' domain knowledge. For example, users inferred that ``{\em we employed 29,239 people...}'' referred to Amazon without the earlier ``Amazon'' mention \sys includes---context needed by LLMs but not humans.}} 


%\agp{For example...}}

\rtwo{In the no-provenance setting, or when the provenance was insufficient (e.g., for RAG), participants often reverted to keyword search. As P7 noted:
“{\em Evidence containing keywords that match the question and answer feels more trustworthy. But I'm not sure if I selected the right keywords due to false positives.}” This keyword search requires multiple passes, typically around 2 distinct keyword search queries, skimming over 18 matches on average
to locate relevant content. 
This inefficiency worsens when keywords 
are frequently present in the document or when documents are large. }
% We also found that participants resorted
% to keyword search when the provenance
% was insufficient (e.g., for RAG).}  

\rtwo{\delete{We found two potential areas of improvement.
First, when the returned provenance was large,
even if minimal,
it was difficult
to interpret.
Providing mechanisms to summarize this provenance
\techreport{(as in relational provenance summarization~\cite{lee2020approximate,alomeir2023summarizing})} may be helpful, such as
highlighting keywords that match the question or answer.}\cradd{We found two areas of improvement. First, large provenance, even if minimal, was hard to interpret; highlighting keywords matching the question or answer could help.} “{\em Highlighting sentences within a long provenance and indicating their correspondence to each part of the answer helps a lot}”,  as noted by P4. 
Second, providing mechanisms to ``tighten'' or ``expand''
the provenance (e.g., dropping the least relevant portions,
or including the most relevant excluded portions), 
depending on the degree of familiarity
with the domain, may also be helpful.  
} 


% \noindent \rthree{{\bf Findings.} We further report insightful findings from the user study regarding the expected provenance representation and how users' domain knowledge affects their judgment. }

% \rthree{
% \squishlist
% \item When the answer contains multiple parts, users expect the relationship between the corresponding provenance for each part of the answer for better reasoning; 
% \item When a provenance spans multiple pages, users expect all provenance snippets to be grouped together in a more concise presentation, while providing the necessary context to connect pieces of information. 
% \item 
% \squishend
% }
